% Generated draft: requires mathematical and source review.
\documentclass[11pt]{article}
\usepackage[margin=25mm]{geometry}
\usepackage{fontspec}
\setmainfont{Latin Modern Roman}
\usepackage{amsmath,amssymb,mathtools,unicode-math}
\setmathfont{Latin Modern Math}
\usepackage{enumitem,needspace,xurl,hyperref}
\setlength{\parindent}{0pt}
\setlength{\parskip}{5pt}
\setlength{\emergencystretch}{3em}
\newcommand{\N}{\mathbb{N}}
\newcommand{\Z}{\mathbb{Z}}
\newcommand{\Q}{\mathbb{Q}}
\newcommand{\R}{\mathbb{R}}
\newcommand{\C}{\mathbb{C}}
\newcommand{\F}{\mathbb{F}}
\newcommand{\A}{\mathbb{A}}
\newcommand{\PP}{\mathbb{P}}
\newcommand{\E}{\mathbb{E}}
\newcommand{\eps}{\varepsilon}
\newcommand{\dd}{\,\mathrm{d}}
\begin{document}
\section*{UnsolvedMath ID: 701775}
\textbf{Source problem code:} catalogue entry
\subsection{Definitions and mathematical statement}
Let \(p\) be a prime. A pro-\(p\) group is a topological group isomorphic to an inverse limit of finite \(p\)-groups. Write \(\Z_p\) for the ring of \(p\)-adic integers and \(\Lambda_G=\Z_p[[G]]\) for the completed group algebra.

A pro-\(p\) Poincar\'e duality group of dimension \(3\) has \(p\)-cohomological dimension \(3\), its trivial module \(\Z_p\) admits a finite resolution by finitely generated projective profinite left \(\Lambda_G\)-modules, and
\[
\operatorname{Ext}^{i}_{\Lambda_G,\mathrm{cont}}(\Z_p,\Lambda_G)
\cong
\begin{cases}
\Z_p,&i=3,\\
0,&i\ne 3,
\end{cases}
\qquad (i\geq 0),
\]
where these Ext groups are computed in the profinite module category and the displayed isomorphism in degree \(3\) is an isomorphism of \(\Z_p\)-modules. No triviality of the natural right \(G\)-action on that degree-three module is imposed. Here \(p\)-cohomological dimension is defined using continuous cohomology with discrete \(p\)-primary torsion modules.

A pro-\(p\) group \(H\) is topologically finitely generated if some finite subset generates a dense subgroup of \(H\). It is finitely presented in the pro-\(p\) category if there are nonnegative integers \(d,r\) and elements \(s_1,\ldots,s_r\) of the free pro-\(p\) group \(F_d\) such that
\[
H\cong F_d/\overline{\langle\!\langle s_1,\ldots,s_r\rangle\!\rangle},
\]
where the denominator is the smallest closed normal subgroup containing those elements. A pro-\(p\) group is coherent if each of its topologically finitely generated closed subgroups is finitely presented in this sense.

\textbf{Question.} For every prime \(p\) and every pro-\(p\) Poincar\'e duality group \(G\) of dimension \(3\), is every topologically finitely generated closed subgroup \(H\leq G\) finitely presented in the pro-\(p\) category?

The supplied source explicitly specifies coherence in the pro-\(p\) sense. It also describes the analogous question for abstract groups as a famous problem and reports an affirmative answer for \(3\)-manifold groups. Those are supplied background statements; the record gives neither precise hypotheses nor a separate reference for the manifold assertion.
\subsection{Short English statement}
Must every \(3\)-dimensional pro-\(p\) Poincar\'e duality group be coherent: does each closed subgroup generated topologically by finitely many elements have a finite pro-\(p\) presentation? The supplied record labels this question open.
\subsection{Sources}
\begin{itemize}
\item Supplied local UnsolvedMath registry record: permanent ID \(701775\), exact source problem code \texttt{catalogue entry}, title \emph{Kourovka Notebook Problem 20.117}. Upstream revision: \texttt{d371ee917ac9647fbb06cd6fc529deec6f3dda82}. The supplied status is \texttt{open}; this status has not been independently verified.
\item \emph{The Kourovka Notebook (Unsolved Problems in Group Theory)}, problem 20.117, attributed to P. Zalesskii, proposed in 2022. The supplied background identifies the 20th issue (2022) and the 21st issue edition (2026), with identifier \texttt{arXiv:1401.0300v48}. These bibliographic details are reproduced from the record without independent source inspection.
\item The supplied collection describes \emph{Kourovka Notebook (Issues 1--21)} as covering 1965--2026 and names E. I. Khukhro and V. D. Mazurov as editors.
\end{itemize}
\end{document}
